\section{Định hướng thiết kế}

Chương này trình bày chi tiết về các lựa chọn công nghệ, kiến trúc hệ thống và các công cụ bổ trợ được nhóm sử dụng để phát triển website doanh nghiệp. Các định hướng này nhằm đáp ứng yêu cầu của giảng viên về việc sử dụng công nghệ thuần túy, đồng thời đảm bảo tính bền vững, hiệu năng và bảo mật cho hệ thống.

\subsection{Mô hình Server-Client}
Website được xây dựng dựa trên kiến trúc Máy chủ - Máy khách (Server-Client). Đây là mô hình tiêu chuẩn giúp tối ưu hóa việc quản lý tài nguyên tập trung và đảm bảo tính nhất quán của dữ liệu:
\begin{itemize}
    \item \textbf{Server (Máy chủ):} Sử dụng máy chủ web Apache kết hợp với trình thông dịch PHP. Máy chủ chịu trách nhiệm tiếp nhận yêu cầu từ người dùng, thực thi các logic nghiệp vụ (business logic), truy vấn cơ sở dữ liệu MySQL và tạo ra các trang HTML động để gửi về máy khách.
    \item \textbf{Client (Máy khách):} Là trình duyệt web của người dùng. Client thực hiện gửi yêu cầu thông qua giao diện người dùng và nhận về kết quả để hiển thị thông tin bằng mã HTML, CSS và các kịch bản JavaScript.
\end{itemize}

\subsection{Kiến trúc MVC (Model-View-Controller)}
Nhằm đảm bảo mã nguồn dễ quản lý và bảo trì, nhóm tự triển khai hệ thống theo mô hình MVC mà không sử dụng framework bên thứ ba. Mô hình này giúp tách biệt rõ ràng các thành phần của ứng dụng:
\begin{itemize}
    \item \textbf{Model (Mô hình dữ liệu):} Chứa các lớp xử lý tương tác trực tiếp với cơ sở dữ liệu MySQL. Thành phần này đảm nhiệm việc truy vấn, thêm mới hoặc cập nhật dữ liệu (như danh sách Q\&A, thông tin liên hệ).
    \item \textbf{View (Giao diện):} Chứa mã nguồn HTML/CSS chịu trách nhiệm hiển thị dữ liệu tới người dùng. Các View được thiết kế để tách rời khỏi logic xử lý PHP giúp giao diện dễ dàng tùy chỉnh.
    \item \textbf{Controller (Bộ điều khiển):} Đóng vai trò trung gian, tiếp nhận yêu cầu từ người dùng, điều phối dữ liệu giữa Model và View để trả về kết quả cuối cùng.
\end{itemize}

\subsection{Ngôn ngữ lập trình và Cơ sở dữ liệu}
\begin{itemize}
    \item \textbf{PHP 7.0+ (Native):} Ngôn ngữ xử lý phía máy chủ chính. Việc sử dụng PHP thuần giúp nhóm kiểm soát hoàn toàn luồng dữ liệu, từ việc xử lý Session cho đến quản lý tệp tin tải lên server.
    \item \textbf{MySQL:} Hệ quản trị cơ sở dữ liệu quan hệ được dùng để lưu trữ toàn bộ thông tin hệ thống. Cấu trúc các bảng dữ liệu được thiết kế tối ưu với các ràng buộc khóa ngoại để đảm bảo tính toàn vẹn dữ liệu.
    \item \textbf{Frontend Core:} Sử dụng HTML5 để định nghĩa cấu trúc, CSS3 (kết hợp thư viện Bootstrap) để thiết kế giao diện Responsive và JavaScript thuần để xử lý các tương tác phía người dùng.
\end{itemize}

\subsection{Giao diện quản trị với Template Srtdash}
Để xây dựng một hệ quản trị nội dung (CMS) chuyên nghiệp, nhóm tích hợp template \textbf{Srtdash}. Đây là lựa chọn chiến lược giúp:
\begin{itemize}
    \item \textbf{Chuyên nghiệp hóa Dashboard:} Cung cấp sẵn các thành phần giao diện quản trị như bảng dữ liệu (DataTables), biểu đồ thống kê và các mẫu form nhập liệu chuẩn hóa.
    \item \textbf{Tính đồng bộ cao:} Đảm bảo giao diện quản lý của Task 1 (Liên hệ/Thông tin chung) và Task 2 (Giới thiệu/Q\&A) có sự nhất quán về trải nghiệm người dùng.
    \item \textbf{Tối ưu hóa thời gian:} Tập trung nguồn lực vào việc xử lý logic PHP backend thay vì phải thiết kế lại giao diện quản trị từ đầu.
\end{itemize}

\subsection{Công cụ quản lý mã nguồn Git và GitHub}
Dự án được quản lý tập trung trên GitHub, giúp nhóm kiểm soát quá trình phát triển một cách chặt chẽ:
\begin{itemize}
    \item \textbf{Quản lý phiên bản:} Theo dõi lịch sử thay đổi của từng dòng code, cho phép nhóm khôi phục lại các trạng thái ổn định nếu xảy ra lỗi trong quá trình thực hiện mô hình MVC.
    \item \textbf{Hợp tác nhóm:} Hỗ trợ việc gộp mã nguồn từ các thành viên làm các nhiệm vụ khác nhau, giúp phát hiện và giải quyết các xung đột mã nguồn (conflicts) kịp thời.
\end{itemize}

\subsection{Định hướng bảo mật và SEO}
\begin{itemize}
    \item \textbf{Bảo mật web:} Triển khai các cơ chế phòng chống lỗ hổng phổ biến như SQL Injection thông qua Prepared Statements, chống XSS bằng cách lọc dữ liệu đầu ra và bảo mật phiên đăng nhập của Admin.
    \item \textbf{Tối ưu hóa SEO:} Cấu trúc mã nguồn thân thiện với bộ máy tìm kiếm bằng cách sử dụng các thẻ Semantic HTML, tối ưu tốc độ tải trang và đảm bảo tính hiển thị tốt trên thiết bị di động (Mobile-friendly).
\end{itemize}

\newpage
\section{Kiến trúc mã nguồn và tổ chức thư mục}

Hệ thống được xây dựng trên ngôn ngữ PHP thuần theo mô hình kiến trúc Model-View-Controller (MVC). Cấu trúc thư mục được phân cấp rõ ràng nhằm tách biệt luồng xử lý dữ liệu, giao diện người dùng và các cấu hình cốt lõi của hệ thống.

\subsection{Sơ đồ cấu trúc thư mục tổng quát}
Dưới đây là sơ đồ phân cấp các thành phần chính trong thư mục gốc \texttt{WEB252}:

\begin{figure}[H]
    \centering
    \includegraphics[width=0.8\textwidth]{image/folder_structure_diagram.png}
    \caption{Sơ đồ phân cấp thư mục dự án theo mô hình MVC}
\end{figure}

\subsection{Chi tiết các thành phần cốt lõi}

Dựa trên cấu trúc thực tế của dự án, các thành phần được phân bổ nhiệm vụ như sau:

\begin{itemize}
    \item \textbf{Thư mục \texttt{app/core/}:} Chứa các thành phần "xương sống" của hệ thống.
    \begin{itemize}
        \item \texttt{Router.php}: Điều hướng các yêu cầu HTTP đến đúng Controller.
        \item \texttt{Database.php}: Quản lý kết nối cơ sở dữ liệu MySQL bằng PDO.
        \item \texttt{BaseController.php \& BaseModel.php}: Các lớp cha cung cấp phương thức dùng chung cho toàn bộ hệ thống.
        \item \texttt{AuthMiddleWare.php}: Kiểm soát quyền truy cập (Access Control) giữa Admin và Client.
    \end{itemize}
    
    \item \textbf{Thư mục \texttt{app/controllers/}:} Chia làm hai phân hệ rõ rệt.
    \begin{itemize}
        \item \texttt{admin/}: Chứa các Controller xử lý nghiệp vụ quản trị (Quản lý Q\&A, Contact, Sản phẩm).
        \item \texttt{client/}: Xử lý các yêu cầu từ phía người dùng cuối như gửi form liên hệ, xem danh sách FAQ.
    \end{itemize}

    \item \textbf{Thư mục \texttt{app/models/}:} Chứa các lớp tương tác trực tiếp với cơ sở dữ liệu. Ví dụ: \texttt{FaqModel.php} cho Task 2 và \texttt{ContactModel.php} cho Task 1.

    \item \textbf{Thư mục \texttt{views/}:} Chứa giao diện hệ thống, sử dụng Template Engine tự viết để tách biệt mã PHP và HTML.
    \begin{itemize}
        \item \texttt{admin/partials/}: Chứa các thành phần giao diện lặp lại như Header, Footer và Sidebar của template \textbf{Srtdash}.
    \end{itemize}

    \item \textbf{Thư mục \texttt{uploads/}:} Lưu trữ tài nguyên tĩnh (hình ảnh sản phẩm, slider, avatar) giúp hệ thống quản lý dữ liệu đa phương tiện hiệu quả.
\end{itemize}

\subsection{Luồng xử lý yêu cầu (Request Lifecycle)}
Mọi yêu cầu từ trình duyệt đều được tập trung về file \texttt{index.php} (Front Controller). Tại đây, hệ thống sẽ thực hiện các bước sau:
\begin{enumerate}
    \item Nạp cấu hình từ \texttt{config/database.php}.
    \item Khởi tạo \texttt{Router} để phân tích URL.
    \item Chạy \texttt{AuthMiddleWare} để kiểm tra quyền hạn.
    \item Gọi đúng \texttt{Controller} và \texttt{Action} tương ứng để trả về kết quả cho người dùng.
\end{enumerate}